\chapter{Structure entière des lecteurs de constantes}
\label{chap:canales-enteros-constantes}
\index{constantes mathématiques!canaux entiers}
\index{racine numérique!résidu et retenue}
\index{demi-couronne}

Les générations précédentes distinguent la fonction de chaque constante. Ce
bloc démontre la structure entière de leurs lecteurs. Les chiffres
\(729,271,315\) et \(161\) sont construits comme \emph{canaux entiers relatifs
à des règles typées déclarées} ; ils n'entrent pas comme approximations décimales. Leurs
statuts ne sont pas identiques : \(729\) est la puissance quadratique du canal
\(27\) ; \(271\) et \(315\) reçoivent une réalisation nonadique construite à partir de
la partition de phase TPK publiée et une seconde réalisation relative à l'élévation de
Steiner ; et
\(161\) reste, dans la portée présente, une identité entière sans
réalisation ensembliste propre. Le catalogue est ensuite classé au moyen de
prédicats et de calibres publiés. Ce n'est qu'à la fin qu'une carte
archimédienne peut agir pour produire des développements décimaux.

Cet ordre sépare deux opérations distinctes. Chercher l'entier \(271\)
parce qu'il rappelle les premiers chiffres de \(e\) serait une lecture
rétrospective. Trouver une région du TPK et l'appeler \(e\) sans construire le
type de retour qu'elle représente serait une dénomination sans application. C'est pourquoi on
sépare les canaux entiers, la classification orbitale conditionnée et le
dictionnaire qui les relie. Le canal entier isolé ne sélectionne pas une
constante ; la sélection appartient au lecteur TPK enrichi, qui incorpore
l'orbite, l'orientation, la survie et la mémoire de frontière sans
recevoir en entrée la valeur décimale qu'il évalue finalement.

\section{Résidu visible, racine numérique et mémoire de couronne}

Pour \(n>0\), la racine numérique nonadique est la section
\[
 \operatorname{dr}_9(n)
 =\begin{cases}
 9,&n\equiv0\pmod9,\\
 n\bmod9,&n\not\equiv0\pmod9.
 \end{cases}
\]
La réduction modulo neuf conserve la classe visible ; la section par
\(\{1,\ldots,9\}\) conserve en outre le choix du représentant nul.
Nous définissons le quotient relevé
\[
 q_9^\sharp(n):=\frac{n-\operatorname{dr}_9(n)}9
 \in\mathbb Z_{\geq0}.
\]
On obtient la décomposition exacte, qui n'est pas la convention usuelle de reste
dans \(\{0,\ldots,8\}\),
\[
 n=9q_9^\sharp(n)+\operatorname{dr}_9(n).
\]
Elle sépare deux lectures : \(\operatorname{dr}_9(n)\) enregistre la valeur
visible et \(q_9^\sharp(n)\) enregistre combien de couronnes complètes précèdent le
représentant visible choisi. C'est la forme exacte de l'image
familière selon laquelle le module compacte et la retenue conserve ce que la
compaction perdrait.

Dans APP, les feuillets additif et multiplicatif produisent
\[
 \sum\Sigma=405,
 \qquad
 \sum\Pi=459,
 \qquad
 \Delta:=\sum\Pi-\sum\Sigma=54.
\]
La somme des quotients additifs est
\[
 S:=\sum Q^+=45.
\]
On définit
\[
 \kappa=\frac{\Delta}{2}=27,
 \qquad
 \Theta=5\kappa=135.
\]
Tous ces nombres proviennent de l'orthogramme et de ses deux évaluations. Aucune
constante n'a encore été utilisée.

Le bloc central
\[
 B_c=\begin{pmatrix}7&2\\2&7\end{pmatrix}
\]
produit quatre mots ordonnés :
\[
 72,\qquad27,\qquad77,\qquad22.
\]
Elles satisfont
\[
 72+27=77+22=99,
 \qquad
 72-27=45,
 \qquad
 77-22=55=54+1.
\]
Le centre ne reste pas isolé : ses deux lectures complémentaires se
projettent sur le bord \(99\), tandis que leurs différences récupèrent la
couronne additive \(45\) et le défaut pointé \(54+1\).

\section{Théorème des quatre canaux entiers}

\begin{theorem}[Identités typées des quatre canaux]
\label{thm:derivacion-canales-enteros}
Relativement aux règles de composition
\(\kappa\mapsto\kappa^2\), \(\Theta\mapsto2\Theta+1\),
\((\Theta,S)\mapsto2\Theta+S\) et
\((\Theta,\kappa)\mapsto\Theta+\kappa-1\), on définit les quatre canaux
entiers
\[
\begin{aligned}
 C_\alpha&=\kappa^2=729,\\
 C_e&=2\Theta+1=271,\\
 C_\pi&=2\Theta+S=315,\\
 C_\varphi&=\Theta+\kappa-1=161.
\end{aligned}
\]
De plus,
\[
\begin{aligned}
 C_e&=5\Delta+1=10\cdot27+1
      =4\Delta+(\Delta+1),\\
 C_\pi&=5\Delta+S=7S,\\
 C_\varphi&=3\Delta-1
      =4\Delta-(\Delta+1),
\end{aligned}
\]
et les fermetures suivantes sont satisfaites
\[
 \boxed{1000=C_\alpha+C_e=729+271}
\]
et
\[
 \boxed{C_\pi+C_e+C_\varphi-3=744=729+15.}
\]
Le canal de \(\alpha\) possède également la lecture centrale
\(729=10\cdot72+9\), et le canal de \(e\), la lecture
\(271=10\cdot27+1\).
\end{theorem}

\begin{proof}
Substituer \(\Delta=54\), \(S=45\), \(\kappa=27\) et
\(\Theta=135\) donne les quatre formules. Les expressions alternatives sont des
identités entières. La première complétion s'obtient à partir de
\(27^2+2(5\cdot27)+1=729+270+1=1000\) ; la seconde, de
\(315+271+161-3=744=729+15\). Les deux divisions décimales découlent des
mots centraux \((72,27)\) et de leurs restes déjà fixés par la
complétion cubique.
\end{proof}

L'égalité \(1000-3^6=271\) est universelle pour cet espace ambiant cubique. À elle
seule, elle ne distingue pas \(e\) : dans un espace ambiant de cardinal \(B\), construit avec
un alphabet de cardinal \(q\) en dimension \(d\), le complément formel est
\(B-q^d\). L'information supplémentaire de cette section est que \(271\) admet
aussi une réalisation comme retour pointé et la lecture centrale \(27|1\).
Son association à l'orbite nommée \(e\) appartient au dictionnaire typé
ultérieur, non à la soustraction isolée.

\section{Demi-couronne nonadique relative à la règle de phase TPK}

La règle minimale de mouvement de l'émetteur TPK ne traite pas les neuf instants d'une
fenêtre comme une liste interchangeable. Relativement à l'origine de phase
publiée, elle définit
\[
 D_9=\{1,\ldots,9\},
 \qquad
 H_+=\{1,4,7\},
 \qquad
 H_-=\{2,5,8\},
 \qquad
 H_0=\{3,6,9\}.
\]
Les instants de \(H_+\) activent la lecture additive, ceux de \(H_-\) la
lecture multiplicative et ceux de \(H_0\) conservent le canal neutre. Posons
\[
 H_{\rm act}=H_+\sqcup H_-,
 \qquad
 O_{\rm ph}=\{1,\ldots,8\},
 \qquad
 \star_{\rm ph}=9,
\]
et soit \(P_2(H_{\rm act})\) l'ensemble des paires non ordonnées de phases
actives. Nous définissons
\[
\begin{aligned}
 \mathcal F^{\rm ph}_6&=H_{\rm act}\times P_2(H_{\rm act}),\\
 \mathcal F^{\rm ph}_8&=O_{\rm ph}\times P_2(H_{\rm act}),\\
 \Theta_{\rm ph}&=D_9\times P_2(H_{\rm act}).
\end{aligned}
\]

\begin{theorem}[Demi-couronne relative à la phase TPK]
\label{thm:semicorona-fase-tpk}
Relativement à l'origine de fenêtre publiée et au retour terminal
\(9\to1\), on a
\[
 |\mathcal F^{\rm ph}_6|=90,
 \qquad
 |\mathcal F^{\rm ph}_8|=120,
 \qquad
 |\Theta_{\rm ph}|=135.
\]
La frontière orientée, le retour pointé, la frontière latérale neutre et la
fermeture latérale sont, respectivement,
\[
\begin{aligned}
 \partial_{\rm ph}&=\Theta_{\rm ph}\times\{-1,+1\},\\
 E_{\rm ph}&=\partial_{\rm ph}\sqcup\{\infty\},\\
 L_{\rm ph}&=H_0\times P_2(H_{\rm act}),\\
 P_{\rm ph}&=\partial_{\rm ph}\sqcup L_{\rm ph},
\end{aligned}
\]
et satisfont
\[
 \boxed{|\partial_{\rm ph}|=270,\quad
 |E_{\rm ph}|=271,\quad |L_{\rm ph}|=45,\quad
 |P_{\rm ph}|=315.}
\]
La construction est covariante sous un changement d'origine : translater la fenêtre
de \(C_9\) translate simultanément toutes les strates et ne change ni leurs
cardinaux ni leurs incidences.
\end{theorem}

\begin{proof}
La règle de signe donne six phases actives et trois neutres, de sorte que

\[
 |P_2(H_{\rm act})|=\binom62=15.
\]

Multiplier par \(6\), \(8\) et \(9\) produit \(90\), \(120\) et \(135\).
L'orientation double \(135\) ; le retour pointé ajoute un élément ; et les
trois phases neutres apportent \(3\cdot15=45\) positions latérales. La
covariance découle du fait qu'une translation de \(D_9\) conjugue la règle de phase
et transporte chaque produit cartésien par une bijection.
\end{proof}

Ce théorème améliore l'identité de cardinalité \(270+1\) de manière substantielle :
l'ordre \(6\subset8\subset9\) n'est plus choisi parmi \(9!\) ordres possibles,
mais provient de l'horloge de phase TPK et de son retour publié. La
canonicité est relative à l'origine de la fenêtre ; la changer ne modifie pas
l'objet à l'action naturelle de \(C_9\) près.

\section{Demi-couronne relative à l'élévation
hexade--octade et retour pointé}

Fixons la dodécade, l'alignement et une hexade
\(H\in\mathcal H(D)\) de la construction de
\cref{sec:w12-w24-elevacion}. D'après
\cref{thm:elevacion-unica-hexada}, il existe une unique octade \(O_H\) telle que
\(O_H\cap D=H\). Écrivons
\[
 P_2(H)=\{\{a,b\}:a,b\in H,\ a\ne b\},
 \qquad |P_2(H)|=\binom62=15,
\]
et ajoutons à \(O_H\) un état de retour \(\star\) :
\[
 \widehat O_H=O_H\sqcup\{\star\},
 \qquad |\widehat O_H|=9.
\]
On définit
\[
 \Theta_H=\widehat O_H\times P_2(H).
\]
En son sein apparaissent les deux strates combinatoires
\[
 \mathcal F_6=H\times P_2(H),
 \qquad
 \mathcal F_8=O_H\times P_2(H).
\]

\begin{theorem}[Réalisation de la demi-couronne]
\label{thm:realizacion-semicorona}
Les objets précédents satisfont
\[
 |\mathcal F_6|=90,
 \qquad
 |\mathcal F_8|=120,
 \qquad
 |\Theta_H|=135.
\]
La frontière orientée
\[
 \partial_H=\Theta_H\times\{-1,+1\}
\]
possède \(270\) états. L'adjonction d'un point de retour,
\[
 E_H=\partial_H\sqcup\{\infty\},
\]
possède \(271\) états. La frontière latérale de la strate à six points est
\[
 L_H=\Theta_H\setminus\mathcal F_6,
 \qquad |L_H|=45,
\]
et la fermeture orientée latérale
\[
 P_H=\partial_H\sqcup L_H
\]
tiene \(315\) estados.
\end{theorem}

\begin{proof}
Les cardinaux sont
\[
 6\binom62=90,
 \qquad
 8\binom62=120,
 \qquad
 9\binom62=135.
\]
L'orientation double \(135\), l'adjonction d'un point ajoute une unité et
la différence entre \(\Theta_H\) et la strate à six points contient
\((9-6)15=45\) éléments. Il en résulte \(270,271\) et
\(270+45=315\).
\end{proof}

Ici, les deux entiers ont des types ensemblistes différents. \(271\) est le
cardinal d'une frontière orientée à laquelle on a adjoint un point ; ce n'est pas,
par ce seul fait, l'orbite de \(e\). \(315\) est le cardinal de cette
frontière avec les \(45\) positions latérales qui n'appartiennent pas à la
strate à six points. L'inclusion \(H\subset O_H\) utilisée ici est bien
l'élévation de Steiner déjà démontrée, relativement à la dodécade et à
l'alignement fixés. La demi-couronne ajoute ensuite le point \(\star\),
l'orientation et la frontière latérale. Cette construction ne fournit pas encore
l'ordre nonadique induit par TPK ni un opérateur d'entrelacement entre les faces du cube
et les incidences de Steiner ; ce sont des applications supplémentaires.

Les deux demi-couronnes partagent le même profil de cardinalité, mais ne doivent pas
être identifiées pour autant. Un isomorphisme non marqué entre les chaînes
\[
 H_{\rm act}\subset O_{\rm ph}\subset D_9
 \quad\text{et}\quad
 H\subset O_H\subset\widehat O_H
\]
qui préserve le point de retour peut permuter librement les six éléments
de la première strate et les deux éléments ajoutés dans la seconde. Il existe
\[
 \boxed{6!\,2!=1440}
\]
calibres de chaîne. L'opérateur d'entrelacement phase--Steiner doit sélectionner l'un
d'entre eux au moyen d'incidences supplémentaires ; l'égalité des cardinaux ne le
construit pas.

\section{Action diédrique et classification sous prédicats déclarés}

Pour une émission
\(U=(u_1,\ldots,u_6)\in\mathcal U_6\), nous définissons le mot visible
\[
 \Psi(U)=((-u_1)\bmod3,\ldots,(-u_6)\bmod3)
 \in\mathbb F_3^6.
\]
Soient les permutations de coordonnées définies explicitement par
\[
 r\cdot(u_1,u_2,u_3,u_4,u_5,u_6)
 =(u_3,u_1,u_2,u_5,u_6,u_4),
\]
\[
 s\cdot(u_1,u_2,u_3,u_4,u_5,u_6)
 =(u_4,u_5,u_6,u_1,u_2,u_3).
\]
Elles engendrent
un groupe diédral \(D_3\simeq S_3\). Le catalogue vérifie
\[
 r\mathcal U_6=\mathcal U_6,
 \qquad
 s\mathcal U_6=\mathcal U_6,
 \qquad
 \Psi(gU)=g\Psi(U),
\]
et conserve les multiplicités microscopiques.

L'action divise les \(243\) mots visibles en \(38\) orbites de taille
six et cinq orbites de taille trois. Sur chaque mot \(w\), écrivons
\[
 n(w)=|\Psi^{-1}(w)|,
 \qquad
 M(w)=\sum_{\Psi(U)=w}|F^{-1}(U)|.
\]

Nous fixons maintenant trois prédicats publiés. Le prédicat \(P_\pi\) exige
\(n(w)=5\), \(M(w)=1008\), les multiplicités
\((432,144,144,144,144)\) et le recensement tritique \((2,3,1)\). Le prédicat
\(P_e\) exige, sur l'orbite complète, des émissions uniponctuelles de masse
microscopique \(144\), l'union de valeurs 
\(\bigcup_{w\in\mathcal O}\operatorname{diag}_c(w)=\{0,3,6\}\) et le recensement
\((2,3,1)\). Le prédicat \(P_\varphi\) conserve les
trois premières conditions de \(P_e\) et exige un recensement équilibré \((2,2,2)\).
Les calibres d'orientation sont \(\operatorname{diag}_c=6\) et la classe
cardinale \(ES\).

\begin{theorem}[Classification orbitale finie conditionnée]
\label{thm:selectores-orbitales-constantes}
Dans le catalogue reconstruit, les prédicats précédents sélectionnent les
orbites uniques suivantes sans consulter les chiffres d'aucune constante.
\begin{enumerate}
  \item Il existe une unique orbite de six mots pour laquelle, dans chaque
  mot, \(n(w)=5\), \(M(w)=1008\) et les cinq multiplicités sont
  \(432,144,144,144,144\). Tous ses mots ont pour recensement tritique
  \((n_0,n_1,n_2)=(2,3,1)\). C'est l'orbite diédrique du type régional de
  \(\pi\).
  \item Parmi les orbites radicales d'émissions uniponctuelles et de masse
  microscopique \(144\) dont l'union orbitale des valeurs
  \(\operatorname{diag}_c\) est \(\{0,3,6\}\), il existe une unique
  orbite de recensement \((2,3,1)\). C'est l'orbite de retour de \(e\).
  \item Dans le même secteur, il existe une unique orbite équilibrée de recensement
  \((2,2,2)\). C'est l'orbite proportionnelle de \(\varphi\).
\end{enumerate}
L'orientation déclarée par
\(\operatorname{diag}_c=6\) et la classe cardinale \(ES\) sélectionne,
respectivement,
\[
 \boxed{
 w_\pi=010211,
 \qquad
 w_e=201101,
 \qquad
 w_\varphi=121200.
 }
\]
\end{theorem}

\begin{proof}
La reconstruction exhaustive des \(468\) émissions vérifie d'abord la
fermeture de l'action et regroupe les \(243\) mots en orbites. Elle applique ensuite
les prédicats indiqués à toutes celles-ci. Le premier prédicat ne laisse qu'une seule
orbite ; les deux prédicats de fibre radicale uniponctuelle laissent, respectivement, une
orbite de recensement \((2,3,1)\) et une de recensement \((2,2,2)\). Enfin, on
applique le calibre d'orientation à chaque orbite et il reste un unique mot.
Aucun des prédicats du sélecteur ne contient les développements décimaux de
\(\pi,e\) ni de \(\varphi\).
\end{proof}

\subsection{Premier encodeur intrinsèque et sa limite}

Dans la strate des mots visibles qui contiennent exactement deux zéros, il existe
un encodeur qui provient bien de la donnée TPK elle-même :
\[
 c_0(w)=\{j\in\{1,\ldots,6\}:w_j=0\}
 \in P_2(\{1,\ldots,6\}).
\]
Il est \(D_3\)-équivariant car une permutation des coordonnées transporte le
support des zéros. Dans les orbites classées par
\cref{thm:selectores-orbitales-constantes}, ses images sont
\[
\begin{aligned}
 c_0(\mathcal O_\pi)
 &=\{\{1,2\},\{1,3\},\{2,3\},\{4,5\},\{4,6\},\{5,6\}\},\\
 c_0(\mathcal O_e)
 &=\{\{1,6\},\{2,5\},\{3,4\}\},
\end{aligned}
\]
où, dans la deuxième ligne, chaque paire possède deux préimages dans l'orbite.
L'orbite de \(\varphi\) partage le premier type de support avec \(\pi\), mais
s'en distingue par sa fibre uniponctuelle, sa multiplicité et son recensement TRIT\@.

L'application \(c_0\) est le premier échelon de l'encodeur registre--paires ; ce n'est pas la
famille complète \(c_k\). Dans les blocs profonds apparaissent (0,1,3) ou (4)
zéros, de sorte que la même formule ne prend pas ses valeurs dans \(P_2\) à toutes les
profondeurs. De plus, \(\{1,\ldots,6\}\) est ici l'hexade des coordonnées du
mot, non une hexade de Witt déjà identifiée. Prolonger \(c_0\) exige
une règle dépendant de la phase et l'opérateur d'entrelacement d'incidence correspondant.

La microfibre canonique de \(\pi\) reste exactement
\[
 \mathscr R_\pi^{\rm micro}=\Psi^{-1}(010211),
\]
formée de cinq régions \(U_6\) distinctes. Les cinq autres mots de
l'orbite diédrique sont des images de calibre de cette forme orbitale ; ils ne constituent pas
cinq nouvelles microfibres canoniques de \(\pi\). De même, les noms
\(e\) et \(\varphi\) s'appliquent après fixation des prédicats et du calibre
publiés. L'énumération est exhaustive \emph{relativement} à ces
prémisses. Prédicats et calibre sont des composantes déclarées du lecteur TPK
enrichi : ils ne sont pas reconstruits à partir des trois statistiques grossières
qui apparaissent ci-dessous, mais à partir de la généalogie complète de l'émission,
de son orientation, de sa frontière et de sa survie.

La somme des six entiers d'une émission, son résidu modulo neuf et le
nombre d'entrées paires ne suffisent pas pour ce théorème. L'émission orientée
de \(e\) et trois des cinq régions de \(\pi\) partagent exactement
\[
 \sum_j u_j=3316,
 \qquad
 \sum_j u_j\equiv4\pmod9,
 \qquad
 \#\{j:u_j\text{ pair}\}=4.
\]
La distinction exige la structure de la fibre, la multiplicité, le secteur
radical et l'orbite complète. Cette ablation démontre pourquoi la parité et
la racine numérique sont des ingrédients du lecteur, mais ne remplacent pas le
transport TPK\@.

\section{Correspondance typée \texorpdfstring{\(271\to e\) et
\(315\to\pi\)}{271 vers e et 315 vers pi}}

Les
\cref{thm:semicorona-fase-tpk,thm:realizacion-semicorona,thm:selectores-orbitales-constantes}
sont indépendants : le premier construit
la demi-couronne à partir de l'horloge de phase, un autre fournit sa réalisation de
Steiner relative et le troisième classe les orbites sous des prédicats fixés. Pour
composer le premier et le troisième résultat, nous déclarons le
dictionnaire de types que le lecteur doit conserver :
\[
\begin{array}{c|c}
\text{type interne}&\text{type orbital}\\
\hline
\text{retour pointé bidirectionnel}&
\begin{gathered}
\text{fibre radicale à un point avec le recensement}\\
\text{de la branche régionale}
\end{gathered}\\
\text{fermeture de frontière plus latérale}&
\text{orbite de clôture à cinq régions}.
\end{array}
\]

\begin{corollary}[Associations typées et orbitales relatives]
\label{cor:lectores-enteros-relativos}
Relativement au dictionnaire, aux prédicats \(P_e,P_\pi\) et aux
calibres d'orientation précédents, la classification orbitale produit les
associations conditionnées
\[
 \boxed{|E_{\rm ph}|=271\rightsquigarrow \mathcal O_e
        \longrightarrow w_e,}
 \qquad
 \boxed{|P_{\rm ph}|=315\rightsquigarrow \mathcal O_\pi
        \longrightarrow w_\pi.}
\]
Dans le second cas, le représentant orienté \(w_\pi=010211\) détermine la
microfibre canonique de cinq régions \(\Psi^{-1}(010211)\). Les associations
s'obtiennent avant l'évaluation archimédienne et n'utilisent pas les valeurs
décimales des constantes.
\end{corollary}

\begin{proof}
Chaque ligne du dictionnaire possède, d'après le théorème orbital, une unique orbite qui
satisfait le prédicat déclaré. La classification s'effectue d'abord au
niveau de l'orbite diédrique complète ; le calibre orienté choisit ensuite le
représentant indiqué. L'unicité est donc relative au dictionnaire,
aux prédicats et aux calibres fixés.
\end{proof}

Le canal \(161=\Theta+\kappa-1=3\Delta-1\) reste une identité entière.
Bien que \(P_\varphi\) classe une orbite équilibrée et que le calibre publié
sélectionne \(w_\varphi=121200\), cette section ne construit pas d'objet
ensembliste de cardinal \(161\) dont le type induise cette orbite. On
n'affirme donc pas ici de lecteur \(161\to\varphi\).

Le corollaire dépend du dictionnaire de types car les cardinaux, pris
isolément, ne définissent pas de fonction vers le catalogue. Dans le lecteur complet,
ce dictionnaire se réalise au moyen des neuf relations de monodromie de
\cref{chap:extensiones-dinamica-cociente}, avec l'orientation et la
survie régionale. La conséquence importante est positive : une fois
le lecteur structurel fixé, l'évaluation ne consulte aucune approximation
décimale de la constante.

\subsection{Dépendance conjointe des canaux dans le sélecteur de produit}

La factorisation exécutable du catalogue et la paramétrisation de sa signature
additive donnent
\[
 \mathcal U_6\simeq\mathcal S_+\times\mathcal S_\times,
 \qquad
 \mathcal S_+\simeq D_9\times\{\pm1\}.
\]
D'autre part,
\[
 \partial_{\rm ph}
 \simeq(D_9\times\{\pm1\})\times P_2(H_{\rm act}).
\]
Ainsi, une tentative fonctionnelle de descente qui conserve à l'identique la
phase et l'orientation se réduit à l'application
\[
 q_\times:\mathcal S_\times\longrightarrow P_2(H_{\rm act}).
\]
C'est le champ absent du catalogue : les lignes publiées se terminent à la
signature produit \(E_{\rm sig}\) et ne contiennent pas de classe
\(P_2(H_{\rm act})\).

\begin{proposition}[Obstruction au sélecteur produit isolé]
\label{prop:nogo-selector-producto-d3}
Il n'existe pas de fonction totale \(D_3\)-équivariante
\[
 q_\times:\mathcal S_\times\longrightarrow P_2(H_{\rm act})
\]
si \(D_3\) agit sur les (26) signatures produit par les permutations
de coordonnées \(r,s\) publiées et sur \(H_{\rm act}\) par
\[
 a(x)=x+3\pmod9,
 \qquad
 j(x)=-x\pmod9.
\]
\end{proposition}

\begin{proof}
L'action sur \(\mathcal S_\times\) a des orbites de tailles
\[
 6+6+6+3+3+1+1.
\]
Ses deux points globalement fixes sont les signatures
\(000000\) et \(555555\). L'action de \(D_3\) sur les six éléments de
\(H_{\rm act}\) est régulière, et son action induite sur les quinze paires a pour
orbites
\[
 6+3+3+3,
\]
sans point globalement fixe. Une fonction équivariante doit envoyer un point
fixe de la source sur un point fixe de la cible, ce qui est impossible.
\end{proof}

L'obstruction identifie précisément l'information que perd la signature
produit isolée : les secteurs fixes ne peuvent être résolus qu'en conservant la phase,
l'orientation et le registre enrichi. L'encodeur \(c_0\) précédent agit dans
la strate à deux zéros, mais ne peut pas remplacer la famille \(c_k\) par une
seule colonne de \(E_{\rm sig}\).

\section{Parité visible, parité de retenue et frontière physique}

La séparation pair--impair peut se formuler sans dépendre d'une étiquette de
particule. Pour chacun des feuillets
\(\bullet\in\{+,\times\}\), posons
\[
 n_+(i,j)=i+j,
 \qquad
 n_\times(i,j)=ij,
\]
et définissons les deux bits
\[
 \epsilon_R^\bullet(i,j)
 =\operatorname{dr}_9(n_\bullet(i,j))\bmod2,
 \qquad
 \epsilon_Q^\bullet(i,j)
 =q_9^\sharp(n_\bullet(i,j))\bmod2.
\]

\begin{theorem}[Factorisation parité--retenue]
\label{thm:factorizacion-paridad-acarreo}
Pour les \(81\) cellules de chaque feuillet APP,
\[
 \boxed{
 n_\bullet(i,j)\bmod2
 =\epsilon_R^\bullet(i,j)+\epsilon_Q^\bullet(i,j)\pmod2.}
\]
Le recensement exact par paires
\((\epsilon_R,\epsilon_Q)=(00,01,10,11)\) est
\[
 \begin{array}{c|rrrr}
  &00&01&10&11\\
 \hline
 +&16&20&20&25\\
 \times&25&5&20&31
 \end{array}
\]
et, par conséquent, les feuillets additif et multiplicatif possèdent des distributions de
parité--retenue distinctes.
\end{theorem}

\begin{proof}
Réduire modulo deux
\(n=9q_9^\sharp(n)+\operatorname{dr}_9(n)\) donne l'identité, car
\(9\equiv1\pmod2\). Dans le feuillet additif, les multiplicités de \(i+j\) sont
\(1,2,\ldots,9,8,\ldots,1\) ; séparer le représentant et le quotient produit la
première ligne. Le même calcul sur les \(81\) produits \(ij\) produit la
seconde. L'énumération des \(162\) cellules vérifie l'identité avant
d'accumuler les comptages.\qedhere
\end{proof}

Sur \(\ell^2(D_9^2)\), les opérateurs diagonaux
\[
 \Gamma_R^\bullet f(i,j)
 =(-1)^{\epsilon_R^\bullet(i,j)}f(i,j),
 \qquad
 \Gamma_Q^\bullet f(i,j)
 =(-1)^{\epsilon_Q^\bullet(i,j)}f(i,j)
\]
sont des involutions commutatives. La parité non réduite se factorise sous la forme
\[
 \boxed{\Gamma_{\rm raw}^\bullet
 =\Gamma_R^\bullet\Gamma_Q^\bullet.}
\]
C'est la formulation opératorielle précise de l'idée selon laquelle la racine numérique
lit le visible tandis que le quotient conserve l'entraînement des couronnes. La
différence entre les deux lignes de la table est aussi une forme exacte de
l'incompatibilité entre additionner et accumuler le produit.

Chaque involution décompose l'espace en secteurs pairs et impairs et fournit
une graduation \(\mathbb Z/2\mathbb Z\) interne. Cependant, une statistique
de Bose--Einstein ou de Fermi--Dirac est une affirmation sur l'échange de
plusieurs réalisations indiscernables, non sur la parité d'une unique cellule.
Sa dérivation requiert un foncteur d'échange qui représente les groupes
symétriques ou de tresses et qui entrelace la dynamique. La parité HMT construit
le domaine gradué sur lequel ce foncteur devra agir ; elle n'est pas identifiée
ici à sa représentation physique.
